\section{Open-Source Artifact and Transferable Mechanism}
\label{sec:artifact}

The public artifact is organized as a fresh repository named
\code{coherent-svm}. It is not a snapshot of the development tree. The
release contains only the code, configurations, schemas, and documentation
needed to explain and validate the transferable mechanism: one rounded kernel
operator must be used consistently by materialization, caching, local
decisions, and global updates.

\noindent Public repository:
\url{https://github.com/yongqiu-star/coherent-svm}.

\subsection{Mechanism interface}

The central invariant can be expressed independently of a particular GPU
microkernel. Let $Q_p(X_i,X_j)$ denote the declared represented kernel,
including input conversion, dot-product precision, distance rules, and cache
encoding. The same $Q_p$ is used at every state-changing boundary:

\begin{lstlisting}[style=codeblock]
for each outer working set W:
    Kww = materialize(X[W], X[W], operator=Qp)
    delta_alpha = solve_local_qp(Kww, state, representation=Qp)

    for each required global column j in W:
        if cache.contains(j):
            column = cache.load(j)              # already in Qp
        else:
            column = materialize(X, X[j], operator=Qp)
            cache.insert(j, column)
        global_gradient += column * delta_alpha[j]

    select_next_working_set(global_gradient, representation=Qp)
\end{lstlisting}

The key invariant is operator identity: a cache hit, a newly produced column,
the local Hessian, and the gradient update all observe values generated by
the same $Q_p$. Mixed-precision accumulation and scalar control are permitted;
silently changing the kernel operator between producers and consumers is not.

\subsection{Portable reference backend}

The public repository provides a portable CUDA C++ and cuBLAS reference
backend. It uses documented APIs rather than architecture-specific inline
PTX. The release is deliberately small: the backend is in
\path{src/portable_fp16_rbf_backend.cu}; the cuML integration is a versioned
patch; \path{coherent_svm/} contains the runner and provenance code; and
\path{artifact/configs/} contains the experiment contracts. This backend is
intended to validate:

\begin{itemize}
  \item representation coherence across the full state transition;
  \item cache hit/miss semantics and equal-logical-capacity comparisons;
  \item integration of materialization, local solve, and global update; and
  \item normalized timing and held-out-quality outputs.
\end{itemize}

The reference backend is not claimed to reproduce the reported
1.351$\times$ speedup. That result was measured with an aggressively optimized
\code{sm\_120a} backend whose development tree still contains
third-party-derived components without complete file-level redistribution
provenance. Publishing that backend before provenance and license boundaries
are resolved would be misleading. It is therefore excluded from the public
release. The measured performance claim and the portable mechanism artifact
are deliberately reported as two different scopes.

\subsection{Public configurations and normalized manifests}

Stable public configuration identifiers describe the experiment without
encoding a host, user, or private workspace. The central files are listed in
\cref{app:public-configs}. Each run records:

\begin{itemize}
  \item a manifest with the configuration digest, public source revision,
        library digest, logical dataset identifier, and platform class;
  \item per-process measurements with method, shape, device class, solver
        parameters, fit status, iterations, support-vector count, and timing;
  \item summary statistics and explicit completion/finite-value checks; and
  \item for the held-out configuration, accuracy, AUC, prediction
        disagreements, McNemar's exact test, and score correlation.
\end{itemize}

This schema makes the artifact useful even on hardware that cannot reproduce
the optimized backend: the semantic contract and the experiment denominator
remain inspectable and comparable.

The public runner does not automate the independent FP32-objective KKT audit
used in the report. That audit remains a model-difference diagnostic rather
than the convergence definition of the coherent low-precision solver.

\subsection{Transferable lesson}

The artifact treats coherent low precision as an algorithm--system interface,
not as an isolated datatype substitution. The same pattern applies to other
iterative GPU algorithms when cached or materialized values influence future
control decisions: define the approximate operator, make every state-changing
consumer agree on it, and then optimize storage and compute within that
contract. Peak throughput is secondary to preserving a productive and
well-defined state trajectory.
